\section{Quantum Electrodynamics (QED)}

\subsection{Gauge Invariance (Classical EM)}
The electromagnetic field is described by a 4-potential $A_\mu$ with field strength
\begin{equation}
F_{\mu\nu}=\partial_\mu A_\nu-\partial_\nu A_\mu.
\end{equation}
Gauge transformation:
\begin{equation}
A_\mu(x)\to A_\mu'(x)=A_\mu(x)+\partial_\mu\alpha(x).
\end{equation}
Since $F_{\mu\nu}$ is invariant, Maxwell theory can be written as
\begin{equation}
\mathcal{L}_{\mathrm{EM}}=-\frac14 F_{\mu\nu}F^{\mu\nu}.
\end{equation}

\subsection{Inclusion of Matter: Minimal Coupling (QED Lagrangian)}
Start from the free Dirac Lagrangian and impose \emph{local} $U(1)$ invariance:
\begin{equation}
\psi(x)\to e^{ie\alpha(x)}\psi(x).
\end{equation}
Introduce the covariant derivative
\begin{equation}
D_\mu := \partial_\mu + ie A_\mu,
\end{equation}
which transforms covariantly:
\begin{equation}
D_\mu\psi \to e^{ie\alpha(x)}D_\mu\psi.
\end{equation}
The full QED Lagrangian:
\begin{equation}
\mathcal{L}_{\mathrm{QED}}
=
\bar{\psi}(i\gamma^\mu D_\mu - m)\psi
-\frac14 F_{\mu\nu}F^{\mu\nu}.
\end{equation}
Expanding:
\begin{equation}
\mathcal{L}_{\mathrm{QED}}
=
\bar{\psi}(i\gamma^\mu\partial_\mu-m)\psi
-\frac14 F_{\mu\nu}F^{\mu\nu}
-e\,\bar{\psi}\gamma^\mu\psi\,A_\mu.
\end{equation}
The interaction term is:
\begin{equation}
\mathcal{L}_{\mathrm{int}}=-e\,\bar{\psi}\gamma^\mu\psi\,A_\mu.
\end{equation}

\subsection{Conserved Current}
Gauge invariance implies a conserved $U(1)$ current:
\begin{equation}
j^\mu=\bar{\psi}\gamma^\mu\psi,
\qquad
\partial_\mu j^\mu=0,
\end{equation}
and Maxwell's equation becomes
\begin{equation}
\partial_\mu F^{\mu\nu}=e\,j^\nu.
\end{equation}

\subsection{Quantization of the EM Field (Key Issue)}
A massless spin-1 field has gauge redundancy, so naive canonical quantization is subtle.
The physical photon has only \textbf{two} transverse polarizations.

To quantize covariantly, introduce a gauge-fixing term. In Lorenz gauge:
\begin{equation}
\partial_\mu A^\mu=0.
\end{equation}
A standard gauge-fixed Lagrangian (covariant gauges):
\begin{equation}
\mathcal{L}
=
-\frac14 F_{\mu\nu}F^{\mu\nu}
-\frac{1}{2\xi}(\partial_\mu A^\mu)^2.
\end{equation}
Here $\xi$ is the gauge parameter. The most common choice:
\begin{equation}
\xi=1 \quad \text{(Feynman gauge)}.
\end{equation}

\subsection{Lorentz-Invariant Propagators}
\paragraph{Fermion propagator.}
\begin{equation}
S_F(p)=\frac{i(\slashed{p}+m)}{p^2-m^2+i\epsilon}.
\end{equation}

\paragraph{Photon propagator (covariant gauge).}
In momentum space:
\begin{equation}
D_{\mu\nu}(k)
=
\frac{-i}{k^2+i\epsilon}
\left(
\eta_{\mu\nu}-(1-\xi)\frac{k_\mu k_\nu}{k^2}
\right).
\end{equation}
In Feynman gauge ($\xi=1$):
\begin{equation}
D_{\mu\nu}(k)=\frac{-i\eta_{\mu\nu}}{k^2+i\epsilon}.
\end{equation}

\subsection{QED Feynman Rules (Momentum Space)}
From
\begin{equation}
\mathcal{L}_{\mathrm{int}}=-e\,\bar{\psi}\gamma^\mu\psi\,A_\mu,
\end{equation}
the vertex factor is:
\begin{equation}
-ie\gamma^\mu.
\end{equation}

\paragraph{Summary of rules.}
\begin{itemize}
\item Fermion propagator:
\[
\frac{i(\slashed{p}+m)}{p^2-m^2+i\epsilon}.
\]
\item Photon propagator (Feynman gauge):
\[
\frac{-i\eta_{\mu\nu}}{k^2+i\epsilon}.
\]
\item Vertex:
\[
-ie\gamma^\mu.
\]
\item External fermion legs:
\[
u^{(s)}(p),\ \bar{u}^{(s)}(p),\ v^{(s)}(p),\ \bar{v}^{(s)}(p).
\]
\item Momentum conservation at each vertex:
\[
(2\pi)^4\delta^{(4)}\Big(\sum p_{\mathrm{in}}-\sum p_{\mathrm{out}}\Big).
\]
\item Closed fermion loop: extra factor $(-1)$.
\end{itemize}

\subsection{QED Processes (Core Examples)}
\subsubsection{Electron--Muon Scattering: $e^-\mu^-\to e^-\mu^-$}
Tree-level diagram: single photon exchange ($t$-channel).

Invariant amplitude:
\begin{equation}
i\mathcal{M}
=
\bar{u}(p_3)(-ie\gamma^\mu)u(p_1)
\left(\frac{-i\eta_{\mu\nu}}{(p_1-p_3)^2+i\epsilon}\right)
\bar{u}(p_4)(-ie\gamma^\nu)u(p_2).
\end{equation}
So
\begin{equation}
\mathcal{M}
=
(-ie)^2
\left[
\bar{u}(p_3)\gamma^\mu u(p_1)
\right]
\frac{\eta_{\mu\nu}}{q^2}
\left[
\bar{u}(p_4)\gamma^\nu u(p_2)
\right],
\qquad q=p_1-p_3.
\end{equation}

\subsubsection{$e^+e^-\to \mu^+\mu^-$ (Annihilation)}
Tree-level diagram: $s$-channel photon exchange.

Amplitude:
\begin{equation}
i\mathcal{M}
=
\bar{v}(p_2)(-ie\gamma^\mu)u(p_1)
\left(\frac{-i\eta_{\mu\nu}}{(p_1+p_2)^2+i\epsilon}\right)
\bar{u}(p_3)(-ie\gamma^\nu)v(p_4).
\end{equation}

\subsubsection{Compton Scattering: $\gamma e^- \to \gamma e^-$}
Two tree-level diagrams:
\begin{itemize}
\item $s$-channel: intermediate electron after absorbing photon
\item $u$-channel: intermediate electron after emitting photon
\end{itemize}
Key point: both are required by gauge invariance.

\subsubsection{M\o ller and Bhabha Scattering}
\begin{itemize}
\item M\o ller: $e^-e^-\to e^-e^-$ (exchange + identical fermion effects)
\item Bhabha: $e^-e^+\to e^-e^+$ (both $s$ and $t$ channels)
\end{itemize}

\subsection{Gauge Invariance and Ward Identity (Essential Idea)}
Gauge invariance implies relations between amplitudes.
A key practical consequence: physical amplitudes are unchanged under
\begin{equation}
\epsilon_\mu(k)\to \epsilon_\mu(k)+c\,k_\mu
\end{equation}
for external photons (with polarization vector $\epsilon_\mu$).

This leads to the Ward identity (schematically):
\begin{equation}
k_\mu \mathcal{M}^\mu = 0.
\end{equation}
This ensures consistency of the theory and cancellation of unphysical polarizations.

\subsection{Summary (Exam Checklist)}
\begin{itemize}
\item QED Lagrangian:
\[
\mathcal{L}=\bar{\psi}(i\gamma^\mu D_\mu-m)\psi-\frac14F_{\mu\nu}F^{\mu\nu}.
\]
\item Gauge transformation:
\[
A_\mu\to A_\mu+\partial_\mu\alpha,\quad \psi\to e^{ie\alpha}\psi.
\]
\item Interaction vertex:
\[
-ie\gamma^\mu.
\]
\item Propagators:
\[
S_F(p)=\frac{i(\slashed{p}+m)}{p^2-m^2+i\epsilon},
\qquad
D_{\mu\nu}(k)=\frac{-i\eta_{\mu\nu}}{k^2+i\epsilon}\ (\xi=1).
\]
\item Typical tree-level processes: scattering via photon exchange, annihilation, Compton scattering.
\item Ward identity expresses gauge invariance: demonstrate $k_\mu \mathcal{M}^\mu=0$.
\end{itemize}
